% Résumé du mémoire.
% Abstract in French.
%
\chapter*{RÉSUMÉ}\thispagestyle{headings}
\addcontentsline{toc}{compteur}{RÉSUMÉ}

Le sujet de ce mémoire est la compilation des problèmes d'optimisation convexe dans
\texttt{CVXPY}, un langage de modélisation. L'utilisateur écrit son problème sous une forme
proche des mathématiques et le langage le traduit en données numériques acceptées par un
solveur. Le coût de cette compilation tient surtout à la construction des matrices creuses
du problème à partir d'arbres d'expressions symboliques. Ces matrices correspondent à la
jacobienne des contraintes et à la hessienne de l'objectif. Leur construction revient donc à
un calcul de dérivées creuses. Nous proposons une nouvelle implémentation de la canonicalisation
qui est basée sur \texttt{SparseDiffEngine}, une librairie C \emph{open source} qui forme les deux
matrices directement à partir de la structure des expressions. 
Le calcul se déroule en deux
phases. La phase symbolique s'exécute une seule fois par famille de problèmes et détermine
les structures creuses des matrices. La phase numérique remplit les valeurs à chaque résolution. Pour un
problème paramétré, seule la phase numérique est répétée d'une résolution à l'autre. Cette
approche constitue une solution de rechange à la mise en cache déjà offerte par
\texttt{CVXPY} et ne limite pas la manière dont les paramètres interviennent dans le
problème. Sur une suite de problèmes de référence, la première compilation est nettement
plus rapide qu'avec les trois implémentations existantes de \texttt{CVXPY}, et résoudre de
nouveau un problème paramétré prend un temps comparable.
